\section{Discussion and Conclusions}
\label{sec:discussion}
\label{sec:conclusions}

\SHARP\ combines a wide field of view with repeated ecliptic imaging
to obtain asteroid positions and photometry with the Xinglong 60/90~cm
Schmidt telescope. Its observing strategy connects three complementary
scales: the sky area accessible over the year, the displacement measured
between successive exposures, and the time baseline accumulated for each
asteroid over multiple nights. Together, these determine the scientific
information contained in the observations and the targets that can be
selected for further study.

The accumulated footprint describes the range of ecliptic longitudes
and latitudes sampled by the survey. The area receiving three visits
within a night adds a requirement on temporal sampling. A field can
contribute to this area after one complete sequence, even if visits on
other nights are interrupted. The distribution of visits per field and
night therefore complements the cumulative sky coverage when assessing
the observations available for motion measurements. Longer baselines
across nights extend the astrometric arc and provide access to additional
rotational phases.

The spacing between exposures is particularly relevant to the independent
motion search. For the nominal 30~min interval, its angular-rate window
of $3$--$63$~\rateunit\ corresponds to displacements of
$1.50$--$31.50\asec$. Near the slow end, distinguishing a moving source
from a stationary counterpart requires the source positions and static
source treatment to be considered together. The magnitude--rate
distribution in Figure~\ref{fig:retention} places the measured linkage
retention in the context of the detections available to the search.
Its denominator selects objects with at least three detections and
suitable angular motion, so the statistic describes the assembly of
those detections into links. Population studies would additionally
require the detection probability as a function of source brightness,
motion, image quality, and observing geometry.

The calibrated photometry provides a basis for assembling asteroid time
series. The selected stellar sequences test the repeatability of the
aperture measurements over the timescale of a night's observations.
Applying these measurements to rotation studies requires selecting
asteroids with adequate phase coverage and accounting for blending,
changing solar phase angle, and the adopted white-light response
\citep{WarnerHarrisPravec2009}. Measurements separated by several nights
can broaden the phase coverage, while also introducing changes in
observing geometry. The per-object distributions of measurement count,
number of nights, and time span provide a practical starting point for
constructing such samples.

The astrometric residuals characterize the agreement between the measured
positions and catalog predictions within the $1\asec$ association
selection. Their magnitude dependence and nightly offsets distinguish
source-dependent scatter from shifts shared by an observing sequence.
Both matter when combining SHARP measurements with other astrometric
data. The residuals incorporate image-centroid uncertainty, the
reference-frame solution, timing, and ephemeris uncertainty; their
interpretation should follow the measurement selection used to obtain
them. The local image examination of sparsely sampled distant-object
associations further illustrates the value of motion evidence when
assigning catalog identities.

The single-night candidates provide measured positions, brightnesses,
and angular motion over short observed arcs. Their image sequences and
trajectory fits establish coherent motion, while connections across
nights constrain the line-of-sight quantities required for orbit
determination \citep{Kubica2007,Holman2018,Moeyens2021}. The associations
with C/2025~Y1 also illustrate the influence of the reference orbit
catalog: cometary motion remains in the candidate sample when the
attribution catalog contains asteroid orbits. These observations offer
concrete starting points for cross-night association and targeted
follow-up.


The first year provides \yrNexp\ exposures over \yrNnight\ nights,
covering $\yrArea~\degsq$ with a median adjacent-visit interval
of \yrCadence~min. The median catalog model depth is \yrDepth~mag,
and the primary associations contain \yrNAssoc\ measurements
linked to \yrNKnown\ catalog identities. The median radial
O$-$C is $\yrOC\asec$ within the $1\asec$ selection, and the
conditional linkage retention is \yrRetention\%. These quantities
summarize the coverage, sensitivity, and sampling delivered by the
first-year observing program.

The survey thus supplies complementary astrometric and photometric
measurements from repeated wide-field imaging. Its cadence and
per-object sampling connect the first-year observations to studies of
asteroid motion and brightness variation. The characterization presented
here provides the observational basis for selecting those scientific
samples and for interpreting the trajectories recovered by the
single-night search.
